\chapter{Alternative Subroutinen}
\label{kap:alternative}

Die Quellen dieses Kapitels benutzen jeweils eigene Bezeichnungen: \textcite{jansen_solisoba2010}
schreiben $\beta$ für die Kapazität und $n_i$ für die Vielfachheiten,
\textcite{gilmore_gomory1961} nennen die Kapazität eine \emph{Lagerlänge}.
Auch für den zentralen Begriff dieses Kapitels - den Inhalt eines einzelnen Bins, hier
\emph{Konfiguration} genannt und in Definition~\ref{def:konfiguration} festgelegt - führt jede
Quelle ein eigenes Wort.
Bei Jansen und Solis-Oba steht $\mathcal{C}$ für die Menge dieser Objekte, bei Gilmore und Gomory
heißen sie \emph{Schnittmuster}.
Hier wird stattdessen so weit wie möglich die in Kapitel~\ref{chap:grundlagen} festgelegte Notation
verwendet: Instanz $\Gamma$, Kapazität $C$, Binzahlen $\mathrm{FFD}[\Gamma, C]$ und
$\mathrm{OPT}[\Gamma, C]$.
Das macht die drei Subroutinen ohne Übersetzungsschritt vergleichbar und erspart es, denselben
Begriff mehrfach zu definieren.
Zwei Zeichen werden in diesem Kapitel gleichwohl neu belegt: $p_1 > \cdots > p_d$ bezeichnet ab
Abschnitt~\ref{sec:jansen} die \emph{verschiedenen} Aufgabengrößen statt der einzelnen Aufgaben,
und $q$ steht in Definition~\ref{def:konfiguration} für die Anzahl der Konfigurationen, nicht für
die Kapazität einer $q$-Packung aus Kapitel~\ref{chap:performanz}.

\section{MFFD als Subroutine in MULTIFIT}
\label{sec:mffd-subroutine}

\subsection{Der Algorithmus MFFD}
\label{sub:mffd}

\emph{Modified First Fit Decreasing} (MFFD) geht auf \textcite{johnson_garey1985} zurück.
Der Algorithmus löst dasselbe Problem wie FFD - Bin Packing bei fester Kapazität nach
Definition~\ref{def:binpacking} - und erhält dieselbe Eingabe: die Aufgaben einer Instanz
$\Gamma$ und eine Kapazität $C$.
Er klassifiziert die Aufgaben relativ zu $C$ in vier Klassen:
\emph{large} ($p_j > C/2$), \emph{medium} ($C/3 < p_j \le C/2$), \emph{small} ($C/6 < p_j \le C/3$)
und \emph{tiny} ($p_j \le C/6$).
Anschließend durchläuft er fünf Phasen \cite{johnson_garey1985}.

\begin{enumerate}[label=\textit{Phase~\arabic*.}, leftmargin=*, itemsep=3pt]
  \item Für jede large-Aufgabe wird ein eigener Bin eröffnet, in absteigender Reihenfolge der
    Größen. Aufgaben der übrigen Klassen eröffnen hier keinen Bin.
  \item Vorwärtslauf über diese Bins. In jeden Bin wird die größte noch unplatzierte
    medium-Aufgabe gelegt, die dort hineinpasst; passt keine, bleibt der Bin unverändert.
    Jeder Bin erhält so höchstens eine medium-Aufgabe.
  \item Rückwärtslauf über diejenigen Bins, die in Phase~2 keine medium-Aufgabe erhalten haben.
    Passen die beiden kleinsten noch unplatzierten small-Aufgaben nicht gemeinsam in den Bin, so
    bleibt er unverändert; andernfalls werden die kleinste noch unplatzierte small-Aufgabe und
    danach die größte noch passende small-Aufgabe eingefügt.
  \item Erneuter Vorwärtslauf über alle Bins. In jeden Bin wird wiederholt die größte noch
    unplatzierte Aufgabe beliebiger Klasse gelegt, die noch hineinpasst, bis keine mehr passt.
  \item Die danach noch unplatzierten Aufgaben werden mit FFD in neue Bins gepackt.
\end{enumerate}

\noindent
Mit $\mathrm{MFFD}[\Gamma, C]$ wird die Anzahl der insgesamt erzeugten Bins bezeichnet.
\textcite{johnson_garey1985} zeigen
\begin{equation}
  \mathrm{MFFD}[\Gamma, C] \;\le\; \tfrac{71}{60}\,\mathrm{OPT}[\Gamma, C] + \tfrac{31}{6};
  \label{eq:mffd-guete}
\end{equation}
\textcite{yue_zhang1995} verschärfen die additive Konstante auf $1$.

Die beiden Faktoren $71/60$ und $11/9$ sind dabei nicht von derselben Art.
Für FFD ist die Schranke \eqref{eq:ffd-guete} scharf \cite{dosa2007}; das asymptotische
Worst-Case-Verhältnis von FFD \emph{beträgt} also $11/9 \approx 1{,}2222$.
Für MFFD ist $71/60 \approx 1{,}1833$ dagegen nur als obere Schranke bewiesen - ob es
angenommen wird, ist hier nicht belegt.
Gesichert ist damit die Aussage: Das asymptotische Worst-Case-Verhältnis von MFFD ist
\emph{höchstens} $71/60$ und damit echt kleiner als das von FFD.

Die Laufzeit bleibt dabei in derselben Größenordnung wie die von FFD.
Sortieren und Klassifizieren kosten $O(n \log n)$.
In den Phasen~2 bis~4 wird jeder der höchstens $n$ Bins konstant oft betrachtet, und insgesamt
werden höchstens $n$ Aufgaben platziert; hält man die Klassenlisten in sortierten Suchstrukturen,
so kostet jede Suche nach der größten noch passenden Aufgabe $O(\log n)$.
Phase~5 ist ein FFD-Aufruf.
Insgesamt ergibt das $O(n \log n)$.

Beide Schranken beziehen sich auf das Bin-Packing-Problem bei fester Kapazität und sind
asymptotischer Natur: Die additive Konstante fällt erst für große $\mathrm{OPT}[\Gamma, C]$ nicht
mehr ins Gewicht.
Ob sich der kleinere Faktor überträgt, wenn MFFD als Subroutine in MULTIFITs Binärsuche eintritt,
ist damit nicht beantwortet.
Der folgende Abschnitt gibt eine Bedingung an, unter der er vollständig entfällt.

\subsection{Kollaps von MFFD auf FFD}
\label{sub:kollaps}

Der Vorteil von MFFD kann innerhalb von MULTIFIT vollständig entfallen, und zwar aus einem
strukturellen Grund, der sich exakt fassen lässt.
Phase~1 eröffnet für jede large-Aufgabe einen eigenen Bin; die Phasen~2 bis~4 durchlaufen ausschließlich die dabei entstandene Binliste und füllen sie auf.
Phase~5 übergibt alle danach noch unplatzierten Aufgaben an FFD.
Gibt es keine large-Aufgabe, so ist die Binliste nach Phase~1 leer, die Phasen~2 bis~4 sind wirkungslos, und Phase~5 wendet FFD auf die vollständige Aufgabenliste an.
Beide Subroutinen erzeugen dann dieselbe Packung, insbesondere dieselbe Binzahl:
\begin{equation}
  p_{\max} \le C/2 \quad\Longrightarrow\quad
  \mathrm{MFFD}[\Gamma, C] = \mathrm{FFD}[\Gamma, C].
  \label{eq:mffd-gleich-ffd}
\end{equation}
Entscheidend ist, dass MULTIFIT nicht beliebige Kapazitäten testet, sondern nur solche oberhalb der unteren Grenze $c_{\mathrm{low}}^{(0)}$ des Startintervalls~\eqref{eq:startintervall}.
Daraus folgt ein Kriterium, das allein von der Instanz abhängt.

\begin{lemma}[Kollaps von MFFD auf FFD]
\label{lem:kollaps}
Sei $\Gamma$ eine Instanz mit $n$ Aufgaben der Größen $p_1, \ldots, p_n$, sei $m \ge 1$ eine
Maschinenzahl und
\[
  \rho(\Gamma) \;:=\; \frac{\sum_{j=1}^{n} p_j}{m \cdot p_{\max}}.
\]
Betrachtet wird MULTIFIT mit dem Startintervall aus \eqref{eq:startintervall}, dessen untere
Grenze $c_{\mathrm{low}}^{(0)}$ ist.
Ist $\rho(\Gamma) \ge 2$, so liefern MULTIFIT+FFD und MULTIFIT+MFFD für jede Iterationszahl $k$ dieselbe Partition; insbesondere stimmen die Makespans überein.
\end{lemma}

\begin{proof}
Zunächst gilt $c_{\mathrm{low}} \le c_{\mathrm{high}}$ als Invariante über alle Iterationen:
Anfangs ist sie vorausgesetzt, und jeder Schritt ersetzt entweder $c_{\mathrm{high}}$ oder
$c_{\mathrm{low}}$ durch das arithmetische Mittel $\tfrac{1}{2}(c_{\mathrm{low}} + c_{\mathrm{high}})$,
das zwischen beiden Grenzen liegt.
Damit ist jede getestete Kapazität $C$ mindestens so groß wie die zu diesem Zeitpunkt aktuelle
untere Grenze, und da $c_{\mathrm{low}}$ im Verlauf der Suche nur wächst, gilt
$C \ge c_{\mathrm{low}}^{(0)}$ für jede getestete Kapazität.

Aus $\rho(\Gamma) \ge 2$ folgt $2\, p_{\max} \le \sum_j p_j / m \le c_{\mathrm{low}}^{(0)} \le C$, also $p_{\max} \le C/2$ für jede getestete Kapazität $C$.
Nach dem oben beschriebenen Phasenablauf erzeugen beide Subroutinen bei jedem einzelnen Test dieselbe Packung, nach \eqref{eq:mffd-gleich-ffd} insbesondere dieselbe Binzahl.
Die Binärsuche trifft damit in jeder Iteration dieselbe Verzweigungsentscheidung; per Induktion über die $k$ Iterationen stimmen die Intervallgrenzen und die jeweils gespeicherte beste Partition zu jedem Zeitpunkt überein.
\end{proof}

\begin{bemerkung}
Der Beweis nutzt von der oberen Intervallgrenze nichts und gilt daher für jede zulässige
Wahl, etwa die hier verwendete Grenze $c_{\mathrm{high}}^{(0)} = \sum_j p_j/m + p_{\max}$ oder
die ursprüngliche Grenze $\max(2\sum_j p_j/m,\, p_{\max})$ von \textcite{coffman1978}.
Er überträgt sich außerdem auf die Variante der Binärsuche, welche nur ganzzahlige Kapazitäten
testet und in Abschnitt~\ref{sub:hm-aufbau} beschrieben wird:
Sind alle $p_j$ ganzzahlig, so folgt aus $\rho(\Gamma) \ge 2$ auch
$\lfloor \sum_j p_j/m \rfloor \ge \lfloor 2 p_{\max} \rfloor = 2 p_{\max}$,
und alle getesteten Kapazitäten bleiben $\ge 2 p_{\max}$.
\end{bemerkung}

Für gleichverteilte Aufgabengrößen lässt sich das Kriterium in das Verhältnis $n/m$ übersetzen.
Sind die $p_j$ unabhängig gleichverteilt auf $[0, P]$, so ist $\mathbb{E}\!\left[\sum_j p_j\right] = nP/2$ und $p_{\max} \le P$, also
\[
  \mathbb{E}[\rho(\Gamma)] \;\ge\; \frac{nP/2}{mP} \;=\; \frac{n}{2m}.
\]
Ab $n/m \ge 4$ erreicht der Erwartungswert von $\rho(\Gamma)$ damit den Wert $2$.
Daraus folgt allerdings weder, dass $\rho(\Gamma) \ge 2$ für eine einzelne Instanz gilt, noch
dass es wahrscheinlich ist: Bei $n/m = 4$ ist die Schranke gerade $2$, und Instanzen mit wenigen
sehr großen und vielen sehr kleinen Aufgaben verletzen $\rho(\Gamma) \ge 2$ auch bei großem
$n/m$.
Lemma~\ref{lem:kollaps} ist von der Größenverteilung dagegen unabhängig.

Der Austausch von FFD gegen MFFD kann sich also nur auf Instanzen mit $\rho(\Gamma) < 2$
auswirken; für die Subroutinen der folgenden Abschnitte gilt das nicht, da
Lemma~\ref{lem:kollaps} allein die Phasen von MFFD benutzt.

Damit ist zugleich gesagt, was Lemma~\ref{lem:kollaps} \emph{nicht} leistet.
Die Bedingung $\rho(\Gamma) \ge 2$ beschreibt eine Klasse von Instanzen; auf ihr liefern beide
Subroutinen dieselbe Ausgabe, über das Verhalten auf allen übrigen Instanzen folgt daraus nichts.
Eine schärfere obere Schranke für MULTIFIT+MFFD folgt daraus nicht: Auf der Klasse stimmen beide
Zuordnungen überein, dort gilt für MULTIFIT+MFFD also dieselbe Schranke wie für MULTIFIT+FFD,
und außerhalb sagt das Lemma nichts.
Wo die Instanz von \textcite{friesen1984} liegt, ist dagegen für die \emph{untere} Schranke
von Belang: Sie hat $\rho = 858/(13 \cdot 40) = 1{,}65 < 2$ und liegt damit außerhalb der Klasse
(Abschnitt~\ref{sec:friesen}).
Läge sie innerhalb, so erzwänge sie das Verhältnis $13/11$ auch für MULTIFIT+MFFD; so aber bleibt
offen, ob MULTIFIT+MFFD eine kleinere Schranke erfüllt.
Das Lemma grenzt also ein, wo ein Unterschied zwischen den beiden Subroutinen überhaupt möglich
ist.
Wie groß der Anteil der Instanzen mit $\rho(\Gamma) < 2$ auf den Benchmark-Datensätzen ist und wie
sich MFFD dort verhält, wird in Abschnitt~\ref{sec:kollaps-empirisch} untersucht.


\section{OPT+1-Algorithmus von Jansen und Solis-Oba}
\label{sec:jansen}

\textcite{jansen_solisoba2010} geben einen Algorithmus für das Cutting-Stock-Problem an, der bei $d$ verschiedenen Objektlängen - in der Notation dieser Arbeit: $d$ verschiedenen
Aufgabengrößen - eine Packung in höchstens $\mathrm{OPT}+1$ Bins berechnet.
Cutting Stock ist Bin Packing in der \emph{High-Multiplicity}-Darstellung: Statt $n$ einzelner Aufgaben werden $d$ verschiedene Größen $p_1 > p_2 > \cdots > p_d > 0$ mit Vielfachheiten $u_1, \ldots, u_d$ angegeben, wobei $\sum_{i=1}^{d} u_i = n$.
Diese Darstellung wird in diesem Kapitel durchgehend verwendet.
Die Implementierung und die Messungen in Abschnitt~\ref{sec:hm} arbeiten mit ganzzahligen
Kapazitäten; die Aussagen dieses Abschnitts verlangen das nur dort, wo es ausdrücklich
dabeisteht - Lemma~\ref{lem:jansen-dual} etwa kommt mit $C \ge p_{\max}$ aus.

\begin{definition}[Konfiguration]
\label{def:konfiguration}
Eine \emph{Konfiguration} zur Kapazität $C$ ist ein Vektor $a = (a_1, \ldots, a_d) \in \mathbb{Z}_{\ge 0}^{d}$ mit
\[
  \sum_{i=1}^{d} a_i\, p_i \;\le\; C
  \qquad\text{und}\qquad
  a_i \le u_i \ \text{ für alle } i .
\]
Sie beschreibt den Inhalt \emph{eines} zulässigen Bins.
Die Menge aller Konfigurationen wird mit $\mathcal{K}(\Gamma, C)$ bezeichnet, ihre Anzahl mit $q = |\mathcal{K}(\Gamma, C)|$.
\end{definition}

\noindent
Eine Packung ist damit eine Auswahl von Konfigurationen mit Wiederholung, die jede Größe oft genug abdeckt.
Mit $\mathrm{J}[\Gamma, C]$ wird im Folgenden die Anzahl der Bins bezeichnet, die der Algorithmus von Jansen und Solis-Oba bei Kapazität $C$ auf $\Gamma$ erzeugt.

\subsection{Laufzeit}
\label{sub:jansen-laufzeit}

Die Laufzeit beträgt \cite[Theorem~1.1]{jansen_solisoba2010}
\[
  d^{O(d2^d)} \times 2^{O(8^d)} \times O\!\left((\log^2 n + \log C)^3 \log^5 n\right).
\]
Dabei erscheint $\log C$, weil die Kapazität Teil der Eingabegröße ist.
Die Laufzeit ist polynomiell in $\log n$, aber doppelt exponentiell in $d$.

\subsection{Implementierung}
\label{sub:jansen-implementierung}

Die Implementierung folgt den drei Schritten des Algorithmus von Jansen und Solis-Oba:
Partitionierung der Aufgaben in große und kleine, Lösung eines gemischt-ganzzahligen linearen
Programms (MILP) über alle großen Konfigurationen und anschließende Rundung der fraktionalen
Lösung.
Als Eingabe dienen drei Parameter: die Liste der $d$ verschiedenen Größen $p = (p_1, \ldots, p_d)$,
die zugehörigen Vielfachheiten $u = (u_1, \ldots, u_d)$ sowie die ganzzahlige Kapazität $C$.
Anstelle von Lenstras Algorithmus \cite{lenstra1983} wird Gurobi eingesetzt:
Gurobi bietet keine polynomielle Worst-Case-Garantie, ist für die hier auftretenden
MILP-Größen aber praktikabel; die gemessenen Laufzeiten stehen in
Abschnitt~\ref{sub:hm-laufzeit}.

Im Originalalgorithmus wird über alle Teilmengen der Größe $2^d$ der weiter unten in Schritt~1
eingeführten Menge $\mathcal{K}^B$ aller großen Konfigurationen iteriert,
da Lenstras Algorithmus eine feste obere Schranke an die Anzahl ganzzahliger Variablen voraussetzt
\cite{jansen_solisoba2010}.
Dass solche Teilmengen genügen, folgt aus einem Ergebnis von \textcite{eisenbrand_shmonin2006}
(dort zitiert in \cite{jansen_solisoba2010}): Das ganzzahlige Programm über alle Konfigurationen
besitzt eine optimale Lösung, in der höchstens $2^d$ Variablen positiv sind.
Mit Gurobi entfällt die Einschränkung: Es genügt, ein einziges MILP mit allen großen Konfigurationen
als ganzzahlige Variablen aufzustellen.
Die Subset-Enumeration wäre in der Praxis auch gar nicht durchführbar: Mit der Schranke
$|\mathcal{K}^B| \le \binom{3d-1}{d}$, welche Schritt~1 begründet, sind es für $d = 3$ bis zu
$\binom{56}{8} \approx 1{,}4 \cdot 10^{9}$ Teilmengen, für $d = 5$ bis zu etwa $10^{70}$ -
verglichen mit einem einzigen MILP, dessen Zahl ganzzahliger Variablen für festes $d$
konstant beschränkt ist.

\subsubsection{Schritt 1: Partitionierung in große und kleine Aufgaben}

Zunächst wird der Schwellenwert
\[
  \varepsilon = \frac{1}{2d - 1}
\]
berechnet. Aufgaben mit Größe $p_i \ge \varepsilon C$ heißen \emph{groß}, alle übrigen \emph{klein};
$\alpha$ bezeichnet die Anzahl der großen Größen, sodass $p_1, \ldots, p_\alpha$ groß sind.
Für eine Konfiguration $a \in \mathcal{K}(\Gamma, C)$ heißt $a^B = (a_1, \ldots, a_\alpha)$ ihr \emph{großer Anteil};
die Menge aller großen Anteile wird mit $\mathcal{K}^B$ bezeichnet.
Die Wahl von $\varepsilon$ ist nicht willkürlich: Da jede große Aufgabe mindestens
$\varepsilon C = C/(2d-1)$ groß ist, passen höchstens $1/\varepsilon = 2d-1$ große Aufgaben in
einen Bin.
Ein großer Anteil ist damit ein Vektor $a^B \in \mathbb{N}_0^{\alpha}$ mit
$\sum_i a^B_i \le 2d-1$, wovon es höchstens $\binom{3d-1}{d}$ viele gibt.
Das begrenzt $|\mathcal{K}^B|$ auf eine Konstante für festes $d$.

Als triviale obere Schranke für die Binärsuche wird $m^* = n = \sum_{i=1}^{d} u_i$ gesetzt,
da jede Aufgabe in einem eigenen Bin untergebracht werden kann.

\subsubsection{Schritt 2: Binärsuche nach dem kleinsten zulässigen \texorpdfstring{$m$}{m}}

Für jeden Kandidaten $m$ wird geprüft, ob die folgende MILP-Instanz zulässig ist
\cite{jansen_solisoba2010}:
\begin{align*}
	&\sum_{\substack{a \in \mathcal{K}(\Gamma, C) \\ a^B = b}} x_{a} \leq y_b,
	\quad \forall\, b \in \mathcal{K}^B \\
	&\sum_{b \in \mathcal{K}^B} y_b = m  \\
	&\sum_{b \in \mathcal{K}^B} b_i\, y_b \geq u_i, \quad \forall\, i = 1, \ldots, \alpha  \\
	&\sum_{a \in \mathcal{K}(\Gamma, C)} a_i\, x_{a} \geq u_i, \quad \forall\, i = \alpha+1, \ldots, d \\
	&x_{a} \geq 0, \quad y_b \in \mathbb{Z}_{\geq 0} \nonumber
\end{align*}

\noindent Die Variablen und Nebenbedingungen bedeuten:

\begin{itemize}
  \item $y_b \in \mathbb{Z}_{\geq 0}$ ist die Anzahl der Bins, deren großer Anteil die Konfiguration
        $b \in \mathcal{K}^B$ ist.
  \item $x_{a} \geq 0$ ist die fraktionale Anzahl der Bins mit vollständiger Konfiguration $a$.
        $x_{a}$ darf gebrochen sein; die Ganzzahligkeit wird erst in Schritt~3 durch Runden hergestellt.
  \item $a_i$ bzw.\ $b_i$ gibt an, wie viele Aufgaben der Größe $p_i$ in der Konfiguration
        $a$ bzw.\ $b$ enthalten sind.
\end{itemize}

Eine Binärsuche über $\{1, \ldots, m^*\}$ findet das kleinste zulässige $m$ in $O(\log n)$
Gurobi-Aufrufen und speichert es zusammen mit der Lösung $(x^+, y^+)$ als neues $m^*$.

\subsubsection{Schritt 3: Rundung und Konstruktion der Packung}

Die $y^+$-Werte sind ganzzahlig; für jedes $b \in \mathcal{K}^B$ werden $y^+_b$ Bins mit den
entsprechenden großen Aufgaben belegt.
Da das MILP $\geq$-Schranken verwendet, kann $y^+$ einzelne Größen öfter als $u_i$-mal nutzen;
diese Aufgaben werden im Nachhinein vom Algorithmus entfernt.
Die fraktionalen $x^+$-Werte werden abgerundet: $\tilde{x}_{a} = \lfloor x^+_{a} \rfloor$.
Für jede Konfiguration $a$ mit $\tilde{x}_{a} \geq 1$ werden $\tilde{x}_{a}$
Bins mit passendem großem Anteil $a^B$ um die kleinen Aufgaben von $a$ ergänzt.
Da $\tilde{x}_{a} = \lfloor x^+_{a} \rfloor \le x^+_{a}$, folgt aus Nebenbedingung~1 des MILP, dass für jedes $b$ genügend Bins
bereitstehen: $\sum_{a^B = b} \tilde{x}_{a} \le y^+_b$.
Kleine Aufgaben, die nach dem Abrunden nicht abgedeckt sind, verteilen
\textcite{jansen_solisoba2010} zunächst mit Next Fit in nicht absteigender Größe auf die
vorhandenen $m^*$ Bins.
Nach ihrem Rundungsargument bleiben dabei höchstens $2d - 1$ Aufgaben übrig, jede kleiner als
$\varepsilon C$, und diese passen gemeinsam in einen einzigen zusätzlichen Bin; die Packung
verwendet damit insgesamt höchstens $m^* + 1$ Bins \cite[Lemma~3.1]{jansen_solisoba2010}.
Die Implementierung verteilt die Reste an dieser Stelle mit First Fit in absteigender Größe und
legt nur das, was danach noch übrig ist, in den zusätzlichen Bin.
Das Rundungsargument von \textcite{jansen_solisoba2010} ist für Next Fit geführt; mit
$\mathrm{J}[\Gamma, C]$ ist im Folgenden deshalb stets die publizierte Fassung gemeint, für die
die Schranke $\mathrm{OPT}[\Gamma, C] + 1$ bewiesen ist.
Ob sie auch für die abweichende Verteilung gilt, ist hier nicht gezeigt;
Abschnitt~\ref{sub:hm-garantien} prüft sie für die gemessene Implementierung empirisch nach.

\subsection{Theoretische Anwendbarkeit als MULTIFIT-Subroutine}
\label{sub:anwendbarkeit}

Die äußere Binärsuche in MULTIFIT akzeptiert $C$ genau dann, wenn $\mathrm{J}[\Gamma, C] \le m$.
Wie schon $\mathrm{FFD}[\Gamma, \cdot\,]$ (Bemerkung~\ref{bem:nicht-monoton}) kann auch
$\mathrm{J}[\Gamma, \cdot\,]$ nicht monoton fallen; ein Zeuge dafür wird an der Implementierung in
Abschnitt~\ref{sub:hm-garantien} gemessen.
Anders als bei FFD lässt sich die Abweichung hier jedoch beziffern.

\begin{lemma}[Begrenzte Nicht-Monotonie]
\label{lem:nichtmonotonie}
	Für alle $C \in \mathbb{N}$, $\delta \in \mathbb{N}_{\ge 1}$ und alle Instanzen~$\Gamma$ gilt
  \[
    \mathrm{J}[\Gamma, C + \delta] \le \mathrm{J}[\Gamma, C] + 1.
  \]
\end{lemma}

\begin{proof}
  Aus der Minimalität von $\mathrm{OPT}[\Gamma, C]$ folgt $\mathrm{J}[\Gamma, C] \ge \mathrm{OPT}[\Gamma, C]$.
  Da des Weiteren jede Packung bei Kapazität $C$ auch bei Kapazität $C+\delta$ gültig ist, gilt
  $\mathrm{OPT}[\Gamma, C+\delta] \le \mathrm{OPT}[\Gamma, C]$. Damit folgt
  \[
    \mathrm{J}[\Gamma, C+\delta] \;\le\; \mathrm{OPT}[\Gamma, C+\delta] + 1 \;\le\; \mathrm{OPT}[\Gamma, C] + 1 \;\le\; \mathrm{J}[\Gamma, C] + 1. \qedhere
  \]
\end{proof}

\begin{korollar}
\label{kor:verwurf}
  Gilt $\mathrm{J}[\Gamma, C] \ge m+2$, so existiert keine gültige Packung von $\Gamma$ in $m$ Bins der Kapazität $C$.
\end{korollar}

\begin{proof}
  Aus der OPT+1-Garantie von \textcite{jansen_solisoba2010} folgt $\mathrm{J}[\Gamma, C] \le \mathrm{OPT}[\Gamma, C] + 1$,
  also $\mathrm{OPT}[\Gamma, C] \ge \mathrm{J}[\Gamma, C] - 1 \ge m+1$. \qedhere
\end{proof}

Lemma~\ref{lem:nichtmonotonie} schließt einen Anstieg nicht aus, begrenzt ihn aber auf~$1$:
Beim Übergang $C \to C+\delta$ kann höchstens der Fall
$\mathrm{J}[\Gamma, C] < \mathrm{J}[\Gamma, C+\delta]$ mit Differenz~$1$ eintreten.
Der Grund liegt in Schritt~3: Die abgerundeten Werte $\lfloor x^+_{a} \rfloor$ hinterlassen kleine Aufgaben, die keiner Konfiguration zugeordnet wurden.
Passen diese Restobjekte in die verbleibenden Bins,
so wird kein Extrabin benötigt und $\mathrm{J}[\Gamma, C] = \mathrm{OPT}[\Gamma, C]$.
Passen sie nicht, entsteht ein zusätzlicher Bin und $\mathrm{J}[\Gamma, C] = m^*+1$, wobei
$m^* \le \mathrm{OPT}[\Gamma, C]$ ist.
Nicht-Monotonie tritt auf, wenn dieser Übergang bei $C$ und $C+\delta$ unterschiedlich ausfällt,
obwohl $\mathrm{OPT}[\Gamma, C] = \mathrm{OPT}[\Gamma, C+\delta]$: Werden die Restobjekte bei $C$ verteilt, bei $C+\delta$ jedoch nicht,
so gilt $\mathrm{J}[\Gamma, C+\delta] = \mathrm{J}[\Gamma, C]+1 > \mathrm{J}[\Gamma, C]$.

Da $\mathrm{J}[\Gamma, C] = m+1$ zweideutig ist - entweder $\mathrm{OPT}[\Gamma, C] = m+1$ (echte Kapazitätsknappheit)
oder $\mathrm{OPT}[\Gamma, C] = m$ (Rundungsartefakt von Schritt~3) -
behandelt die Binärsuche diesen Fall konservativ als Ablehnung.
Bleibt $\mathrm{J}[\Gamma, C] = m+1$ für jede getestete Kapazität, so liefert die Suche keine
Packung in $m$ Bins; auf den Instanzen aus Abschnitt~\ref{sub:hm-guete} tritt dieser Fall
nicht auf.
Wird eine Kapazität akzeptiert, so konvergiert die Suche gleichwohl möglicherweise nicht
gegen die kleinste zulässige.

\paragraph{Güte und Einschränkungen.}
Da $\mathrm{J}[\Gamma, C] \le m$ stets $\mathrm{OPT}[\Gamma, C] \le m$ impliziert - nach
Definition des Optimums gilt $\mathrm{OPT}[\Gamma, C] \le \mathrm{J}[\Gamma, C]$ -,
akzeptiert die Suche nur Kapazitäten, bei denen tatsächlich eine Packung in $m$ Bins existiert.
Weil sie $C$ mit $\mathrm{J}[\Gamma, C] = m+1$ (und potenziell $\mathrm{OPT}[\Gamma, C] = m$) verwirft,
konvergiert sie im schlechtesten Fall gegen den optimalen Makespan auf $m-1$ Maschinen.
Mit der Kurzschreibweise $C^{*}_{\max}(\Gamma) = C^{*}_{\max}(\Gamma, m)$ aus
Definition~\ref{def:instanz} trägt folgende Kette:
Für $C \ge C^{*}_{\max}(\Gamma, m-1)$ ist $\mathrm{OPT}[\Gamma, C] \le m-1$ nach
Lemma~\ref{lem:dualitaet}, also $\mathrm{J}[\Gamma, C] \le m$ nach der $\mathrm{OPT}+1$-Garantie
- solche Kapazitäten akzeptiert die Suche stets.
Liegt $C^{*}_{\max}(\Gamma, m-1)$ im Startintervall, so ist der im Grenzwert $k \to \infty$
zurückgegebene Makespan damit durch $C^{*}_{\max}(\Gamma, m-1)$ beschränkt, das Verhältnis zum
Optimum also durch $C^{*}_{\max}(\Gamma, m-1)/C^{*}_{\max}(\Gamma, m)$; bei endlichem $k$ kommt
der Term $2^{-k}$ hinzu.
Das ist eine Schranke je Instanz; als Güte im Sinn von Definition~\ref{def:guete} gelesen, also
als Supremum über alle Instanzen, liefert sie nur den Wert~$2$ - siehe das Beispiel am Ende
dieses Abschnitts - und damit eine weit schwächere Aussage als $13/11$.
Bestimmt in beiden Fällen die Durchschnittslast das Optimum, gilt also
$C^{*}_{\max}(\Gamma, m) \approx \sum_j p_j/m$ und
$C^{*}_{\max}(\Gamma, m-1) \approx \sum_j p_j/(m-1)$, so ist
\[
  \frac{C^{*}_{\max}(\Gamma, m-1)}{C^{*}_{\max}(\Gamma, m)}
  \;\approx\; \frac{\sum_j p_j/(m-1)}{\sum_j p_j/m} \;=\; \frac{m}{m-1},
\]
und $m/(m-1) < 13/11$ gilt für alle $m \ge 7$.
Das ist eine Abschätzung unter dieser Annahme und keine Schranke: Sie setzt voraus, dass beide
Optima von der Durchschnittslast bestimmt werden, und trägt nicht, sobald einzelne große Aufgaben
sie nach oben drücken.
Im Allgemeinen kann $C^{*}_{\max}(\Gamma, m-1)/C^{*}_{\max}(\Gamma, m)$ bis zu $2$ betragen:
$m$ Aufgaben der Größe~$1$ auf $m$ Maschinen liefern $C^{*}_{\max}(\Gamma, m) = 1$ und $C^{*}_{\max}(\Gamma, m-1) = 2$.

Der Algorithmus setzt schließlich zwingend voraus, dass $d$ konstant ist.
Nur dann ist die Laufzeit polynomiell in $\log n$; für allgemeines Scheduling mit $d = n$ macht
die doppelt exponentielle Abhängigkeit den Algorithmus praktisch unbrauchbar.
Anwendbar bleibt das Verfahren daher nur für High-Multiplicity-Instanzen mit kleinem $d$;
Abschnitt~\ref{sub:hm-laufzeit} misst den Bereich $d \le 6$.

\subsection{Modifikation zu einem \texorpdfstring{$\varepsilon$}{ε}-dualen Approximationsalgorithmus}
\label{sub:eps_dual}

\subsubsection{Primale und duale Approximation}

Approximationsalgorithmen für Bin Packing können die Gütegarantie auf zwei Weisen erfüllen.
Der \emph{primale} Ansatz hält die Kapazität~$C$ ein und verwendet dafür bis zu
$(1+\varepsilon)\cdot\mathrm{OPT}[\Gamma, C]$ Bins.
Der \emph{duale} Ansatz kehrt diese Priorität um: Er verwendet höchstens
$\mathrm{OPT}[\Gamma, C]$ Bins und lässt dafür eine Kapazitätsüberschreitung um den
Faktor $(1+\varepsilon)$ zu.

\begin{definition}[$\varepsilon$-dualer Approximationsalgorithmus {\normalfont\cite{hochbaum_shmoys_1987}}]
\label{def:eps-dual}
  Ein Polynomialzeitalgorithmus~$\mathcal{A}$ heißt \emph{$\varepsilon$-dualer Approxi\-mationsalgorithmus}
  für Bin Packing, wenn er für jede Instanz~$\Gamma$ und jede Kapazität $C \ge p_{\max}$ eine
  Packung berechnet, die
  \begin{enumerate}[label=(\roman*)]
    \item höchstens $\mathrm{OPT}[\Gamma, C]$ Bins verwendet, und
    \item jeden Bin mit Gesamtlast $\le (1+\varepsilon)\, C$ befüllt.
  \end{enumerate}
\end{definition}

Der duale Ansatz ist das entscheidende Werkzeug für das Makespan-Problem.
Der Zusammenhang beider Optima ist genau die Dualität aus Lemma~\ref{lem:dualitaet},
$\mathrm{OPT}[\Gamma, C] \le m \iff C^{*}_{\max}(\Gamma) \le C$;
\textcite{hochbaum_shmoys_1987} formulieren sie über skalierte Aufgabengrößen.
Ein $\varepsilon$-dualer Bin-Packing-Algorithmus lässt sich daher direkt in eine Binärsuche über
die Kapazität einbetten: Sie halbiert das Kapazitätsintervall und übernimmt jede Kapazität, bei
der er höchstens $m$ Bins ausgibt.
Diese Prozedur ist strukturell identisch mit MULTIFITs Binärsuche.

\begin{theorem}[Hochbaum und Shmoys {\normalfont\cite{hochbaum_shmoys_1987}}, Theorem~1]
\label{thm:hs-makespan}
  Die Binärsuche über die Kapazität benutze einen $\varepsilon$-dualen Bin-Packing-Algorithmus als
  Subroutine; ihr Startintervall enthalte $C^{*}_{\max}(\Gamma)$ und seine Breite übersteige
  $C^{*}_{\max}(\Gamma)$ nicht.
  Wird sie nach $k$~Iterationen abgebrochen, so besitzt der erzeugte Schedule einen Makespan von
  höchstens
  \[
    (1 + \varepsilon)(1 + 2^{-k})\cdot C^{*}_{\max}(\Gamma).
  \]
\end{theorem}

Für $k \to \infty$ konvergiert die Schranke gegen $(1+\varepsilon)\cdot C^{*}_{\max}(\Gamma)$,
die Güte also gegen $1+\varepsilon$.
Die Konstruktion eines $\varepsilon$-dualen Bin-Packing-Algorithmus ist damit der entscheidende Baustein.

\subsubsection{Jansen-Dual: Modifikation des Rundungsschritts}

Der Algorithmus aus Abschnitt~\ref{sub:jansen-implementierung} gibt höchstens $\mathrm{OPT}[\Gamma, C]+1$
Bins aus; der Extrabin entsteht in Schritt~3, weil die kleinen Aufgaben, die nach dem Verteilen
auf die vorhandenen Bins übrig bleiben, dort in einen zusätzlichen Bin der Kapazität~$C$ gepackt
werden \cite{jansen_solisoba2010}.
Die folgende Modifikation verteilt diese Aufgaben stattdessen auf bestehende Bins und eliminiert
den Extrabin - auf Kosten einer geringen Kapazitätsüberschreitung.

\medskip\noindent
\textbf{Jansen-Dual (modifizierter Schritt~3).}
Schritte~1 und~2 sind identisch zum ursprünglichen Algorithmus.
Schritt~3 ändert sich wie folgt:
\begin{enumerate}[label=\arabic*.]
  \item Runde ab: $\bar{x}_{a} = \lfloor x^+_{a} \rfloor$ für alle Konfigurationen~$a$.
  \item Weise jedem der $m^*$ Bins seine große Konfiguration und die abgerundeten kleinen
        Aufgaben zu. Wie im Original werden Aufgaben, die dabei öfter als $u_i$-mal abgedeckt
        sind, im Nachhinein entfernt; jede Aufgabe liegt danach in höchstens einem Bin.
  \item Sammle die verbliebenen kleinen Aufgaben $L$ (nicht abgedeckt durch die gerundeten $\bar{x}$-Werte).
        Anders als im Original werden sie nicht vorab mit Next Fit verteilt; der Beweis von
        Lemma~\ref{lem:jansen-dual} braucht keine Schranke an $|L|$.
  \item Platziere jede Aufgabe aus~$L$ in einem Bin mit aktueller Last $< C$.
  \item Gib die $m^*$ Bins aus - kein zusätzlicher Bin wird geöffnet.
\end{enumerate}

\begin{lemma}[Jansen-Dual ist $\varepsilon$-dual für $\varepsilon = 1/(2d-1)$]
\label{lem:jansen-dual}
  Sei $\Gamma$ eine Instanz mit $n$ Aufgaben in $d$~verschiedenen Größen, $C \ge p_{\max}$
  eine Kapazität und $\varepsilon = 1/(2d-1)$; $d$ sei konstant.
  Jansen-Dual berechnet eine Packung aller Aufgaben von~$\Gamma$, für die gilt:
  \begin{enumerate}[label=(\roman*)]
    \item die Anzahl der verwendeten Bins beträgt höchstens $\mathrm{OPT}[\Gamma, C]$;
    \item jeder Bin trägt eine Gesamtlast von höchstens $(1+\varepsilon)\, C$.
  \end{enumerate}
\end{lemma}

\begin{proof}
\textit{(i) Binzahl $\le \mathrm{OPT}[\Gamma, C]$.}\;
$m^*$ ist das kleinste $m$, für das das MILP in Schritt~2 zulässig ist.
Da die optimale Packung in $\mathrm{OPT}[\Gamma, C]$ Bins eine zulässige MILP-Lösung liefert,
ist das MILP für $m = \mathrm{OPT}[\Gamma, C]$ zulässig; daher gilt $m^* \le \mathrm{OPT}[\Gamma, C]$ per Definition.
Jansen-Dual öffnet keinen zusätzlichen Bin und verwendet daher genau $m^* \le \mathrm{OPT}[\Gamma, C]$ Bins.

\textit{(ii) Gesamtlast jedes Bins $\le (1+\varepsilon)\, C$.}\;
Jede Konfiguration~$a$ erfüllt $\sum_i p_i\, a_i \le C$,
weshalb jeder Bin nach dem Abrunden in Schritt~3.1 eine Last $\le C$ trägt.
Die verbleibenden Aufgaben aus~$L$ sind ausschließlich kleine Aufgaben mit Größe $< \varepsilon C$.

Aus der Zulässigkeit des MILPs folgt zunächst, dass die Gesamtgröße aller Aufgaben höchstens
$m^* C$ beträgt.
Für $b \in \mathcal{K}^B$ sei dazu $g_b = \sum_{i \le \alpha} p_i\, b_i$ die Größe der großen
Aufgaben eines Bins mit großem Anteil~$b$; wegen $b \in \mathcal{K}^B$ ist $g_b \le C$.
Die Bedarfszeilen der großen Größen liefern
\[
  \sum_{i \le \alpha} p_i\, u_i \;\le\; \sum_{i \le \alpha} p_i \sum_{b} b_i\, y_b
  \;=\; \sum_{b} g_b\, y_b .
\]
Jede Konfiguration~$a$ hält die Kapazität ein, weshalb ihr kleiner Anteil
$\sum_{i > \alpha} p_i\, a_i \le C - g_{a^B}$ erfüllt.
Mit den Bedarfszeilen der kleinen Größen und der ersten Nebenbedingung folgt
\[
  \sum_{i > \alpha} p_i\, u_i
  \;\le\; \sum_{a} \Bigl( \sum_{i > \alpha} p_i\, a_i \Bigr) x_{a}
  \;\le\; \sum_{b} \bigl( C - g_b \bigr) \sum_{a^B = b} x_{a}
  \;\le\; \sum_{b} \bigl( C - g_b \bigr) y_b .
\]
Addiert man beide Ungleichungen, so heben sich die $g_b$ auf und es bleibt
\[
  \sum_{i=1}^{d} p_i\, u_i \;\le\; C \sum_{b} y_b \;=\; m^* C .
\]
Entscheidend ist dabei, dass große und kleine Aufgaben desselben Bins sich \emph{eine} Kapazität
teilen; schätzt man die beiden Bedarfszeilen unabhängig voneinander ab, so liefern sie nur
$2\, m^* C$ und tragen die folgende Schlussweise nicht.

Angenommen nun, beim Platzieren einer Aufgabe $s \in L$ in Schritt~3.4 wäre kein Bin
mit aktueller Last $< C$ mehr vorhanden, d.\,h.\ alle $m^*$ Bins trügen eine Last von mindestens~$C$.
Dann betrüge die Gesamtlast der Bins mindestens $m^* C$.
Da~$s$ noch nicht platziert ist, ist sie zugleich strikt kleiner als die Gesamtgröße aller
Aufgaben, sodass
\[
  m^* C \;<\; \sum_{i=1}^{d} p_i\, u_i
\]
folgte - im Widerspruch zur soeben gezeigten Ungleichung.

Also findet Schritt~3.4 stets einen Bin mit aktueller Last $< C$;
nach dem Einfügen von~$s$ ist die Last dieses Bins kleiner als $C + \varepsilon C = (1+\varepsilon)\, C$.

\textit{Laufzeit.}\; Die Größe des MILP aus Schritt~2 hängt für festes~$d$ nach
Abschnitt~\ref{sub:jansen-laufzeit} allein von~$d$ ab; Schritt~3.4 platziert dagegen jede
Aufgabe aus~$L$ einzeln und kostet mit einem Zeiger über die Bins $O(|L| + m^{*})$ Schritte mit
$|L| \le n$, sodass die Laufzeit polynomiell in~$n$ ist.
Bei \textcite{jansen_solisoba2010} erlaubt die Schranke $|L| \le 2d-1$ eine von~$n$ unabhängige
Restbehandlung; für Korollar~\ref{kor:multifit-jansen-dual} ist der Unterschied ohne Belang.

Alle Aufgaben von~$\Gamma$ werden gepackt: große und abgerundete kleine Aufgaben durch Schritt~3.2,
die verbliebenen kleinen Aufgaben aus~$L$ durch Schritt~3.4.
\end{proof}

\begin{korollar}[MULTIFIT mit Jansen-Dual]
\label{kor:multifit-jansen-dual}
  MULTIFIT mit Jansen-Dual als Subroutine liefert nach $k$~Iterationen
  der Binärsuche einen Schedule mit Makespan von höchstens
  \[
    (1+\varepsilon)(1+2^{-k})\cdot C^{*}_{\max}(\Gamma),
    \qquad \varepsilon = \frac{1}{2d-1}.
  \]
\end{korollar}

\begin{proof}
  Nach Lemma~\ref{lem:jansen-dual} ist Jansen-Dual ein $\varepsilon$-dualer Approximationsalgorithmus
  für Bin Packing.
  MULTIFITs Startintervall \eqref{eq:startintervall} erfüllt die Voraussetzung von
  Theorem~\ref{thm:hs-makespan}: Es enthält $C^{*}_{\max}(\Gamma)$ nach
  Lemma~\ref{lem:startintervall}(a) und~(c), und für seine Breite gilt
  \[
    c_{\mathrm{high}}^{(0)} - c_{\mathrm{low}}^{(0)}
    \;=\; \min\Bigl( \tfrac{1}{m}\textstyle\sum_j p_j,\; p_{\max} \Bigr)
    \;\le\; C^{*}_{\max}(\Gamma)
  \]
  nach Lemma~\ref{lem:untere-schranken}.
  Die Gütegarantie folgt damit aus Theorem~\ref{thm:hs-makespan}.
\end{proof}

\paragraph{Implementierungsoptimierung: Wegfall der inneren Binärsuche.}
Jansens Schritt~2 sucht das minimale $m^*$ per Binärsuche und braucht dafür $O(\log n)$ MILP-Aufrufe.
Im MULTIFIT-Kontext ist die Maschinenzahl~$m$ jedoch von außen bekannt und fest.
Eine Subroutine wie FFD berechnet (approximativ) eine Packung mit möglichst wenigen Bins, und
MULTIFIT vergleicht deren Binzahl mit~$m$.
Für diesen Vergleich genügt aber schon die Antwort auf die Frage, ob eine Packung in höchstens
$m$ Bins der Kapazität~$C$ gefunden wird.
Diese Frage beantwortet ein \emph{einziger} MILP-Aufruf mit der Nebenbedingung
$\sum_{b} y_b \leq m$:
\begin{itemize}[nosep]
  \item \textbf{Lösbar:} Die Packung wird wie in Schritt~3 konstruiert und zurückgegeben.
  \item \textbf{Unlösbar:} $C$ ist zu klein; MULTIFIT erhöht die Kapazität.
\end{itemize}
Diese Variante minimiert die Binzahl nicht mehr, erfüllt also Eigenschaft~(i) von
Definition~\ref{def:eps-dual} nicht: Sie gibt bis zu~$m$ Bins aus, auch wenn
$\mathrm{OPT}[\Gamma, C] < m$ ist.
Erhalten bleibt genau das, was Theorem~\ref{thm:hs-makespan} für die Binärsuche braucht:
Der Aufruf ist für jedes $C \ge C^{*}_{\max}(\Gamma)$ lösbar, gibt höchstens $m$ Bins aus und
hält die Last~$\le (1+\varepsilon)\,C$.
Für die Lastschranke trägt das Widerspruchsargument aus dem Beweis von
Lemma~\ref{lem:jansen-dual} mit~$m$ an Stelle von~$m^*$, sofern die Lösung alle~$m$ Bins
belegt; belegt sie weniger, ist die Last des am wenigsten gefüllten Bins erst recht $< C$.
Die Anzahl der MILP-Aufrufe reduziert sich dadurch von
$O(k \log n)$ auf $O(k)$, wobei $k$ die Anzahl der MULTIFIT-Iterationen ist.

Diese Verkürzung ist an die Modifikation gebunden und lässt sich nicht auf den unveränderten
$\mathrm{OPT}+1$-Algorithmus übertragen.
Dort entscheidet über die Annahme einer Kapazität der Wert $\mathrm{J}[\Gamma, C]$, und der ist
$m^*$ oder $m^*+1$: Der zusätzliche Bin aus Schritt~3 wird geöffnet, wenn nach dem Abrunden der
$x^+$-Werte kleine Aufgaben übrig bleiben, die sich nicht mehr auf die vorhandenen Bins
verteilen lassen.
Ein einzelner Aufruf mit dem festen $m$ an Stelle von $m^*$ liefert in diesem Fall $m+1$ Bins
und verwirft die Kapazität, obwohl das MILP für $m$ Bins lösbar war.
Die Entscheidung hinge dann nicht mehr an der Lösbarkeit, sondern am Rundungsrest.
Erst die $\varepsilon$-duale Fassung, die diese Reste nach Lemma~\ref{lem:jansen-dual} auf die
bereits geöffneten Bins verteilt, macht den zusätzlichen Bin - und mit ihm die Suche nach
$m^*$ - entbehrlich.
Der Vergleich in Abschnitt~\ref{sec:hm} misst beide Fassungen in genau dieser Form: den
$\mathrm{OPT}+1$-Algorithmus mit innerer Binärsuche, seine $\varepsilon$-duale Variante ohne.

\section{Das Konfigurations-LP von Gilmore und Gomory als Subroutine}
\label{sec:gg-lp}

FFD und MFFD sind kombinatorische Verfahren, welche die Aufgaben nach festen Regeln
nacheinander platzieren; der Algorithmus von Jansen und Solis-Oba
löst dagegen ein ganzzahliges Programm über Konfigurationen und rundet dessen Lösung.
Dieselbe Modellierung liegt schon einem klassischen Zugang zum Bin Packing zugrunde, dem linearen
Programm über Konfigurationen - kurz \emph{Konfigurations-LP} - von
\textcite{gilmore_gomory1961,gilmore_gomory1963}.
Dieser Abschnitt beschreibt das Modell, seine Rundung zu einer zulässigen Packung und
den Einbau des Ergebnisses als Subroutine in MULTIFITs Binärsuche.

\subsection{Das Konfigurationsmodell}
\label{sub:gg-modell}

Grundlage ist dieselbe High-Multiplicity-Darstellung wie in Abschnitt~\ref{sec:jansen}: eine
Instanz $\Gamma$ mit $d$ verschiedenen Größen $p_1 > \cdots > p_d > 0$ und Vielfachheiten
$u_1, \ldots, u_d$, eine feste Kapazität $C$ und die Konfigurationsmenge
$\mathcal{K}(\Gamma, C)$ aus Definition~\ref{def:konfiguration}.
\textcite{gilmore_gomory1961} nennen eine Konfiguration ein \emph{Schnittmuster}.
Mit einer Variablen $x_a$ je Konfiguration - der Anzahl der nach Muster $a$ gepackten Bins -
lautet das ganzzahlige Programm (IP)
\begin{equation}
  \min \sum_{a \in \mathcal{K}(\Gamma, C)} x_a
  \quad\text{u.d.N.}\quad
  \sum_{a \in \mathcal{K}(\Gamma, C)} a_i\, x_a \;\ge\; u_i \ \ (i = 1, \ldots, d),
  \qquad x_a \in \mathbb{Z}_{\ge 0}.
  \tag{IP}
  \label{eq:gg-ip}
\end{equation}
Die $\ge$-Form ist keine Abschwächung: Deckt eine Lösung eine Größe öfter ab als nötig, so
entfernt man die überzähligen Aufgaben aus den betroffenen Bins; das Ergebnis ist wieder eine
Auswahl von Konfigurationen und braucht nicht mehr Bins als zuvor.
Umgekehrt liefert jede Packung in $k$ Bins eine zulässige Lösung von \eqref{eq:gg-ip} mit
Zielfunktionswert $k$.
Der Optimalwert von \eqref{eq:gg-ip} ist daher genau $\mathrm{OPT}[\Gamma, C]$ aus
Definition~\ref{def:binpacking}.

\begin{definition}[Konfigurations-LP]
\label{def:gg-lp}
Ersetzt man in \eqref{eq:gg-ip} die Bedingung $x_a \in \mathbb{Z}_{\ge 0}$ durch $x_a \ge 0$, so
entsteht das \emph{Konfigurations-LP}
\begin{equation}
  \min \sum_{a \in \mathcal{K}(\Gamma, C)} x_a
  \quad\text{u.d.N.}\quad
  \sum_{a \in \mathcal{K}(\Gamma, C)} a_i\, x_a \;\ge\; u_i \ \ (i = 1, \ldots, d),
  \qquad x_a \ge 0 ;
  \tag{LP}
  \label{eq:gg-lp}
\end{equation}
sein Optimalwert wird mit $\mathrm{LP}[\Gamma, C]$ bezeichnet.
\end{definition}

Das zugehörige duale Programm (D) lautet
\begin{equation}
  \max \sum_{i=1}^{d} u_i\, y_i
  \quad\text{u.d.N.}\quad
  \sum_{i=1}^{d} a_i\, y_i \;\le\; 1 \ \ \text{für alle } a \in \mathcal{K}(\Gamma, C),
  \qquad y \ge 0 .
  \tag{D}
  \label{eq:gg-dual}
\end{equation}
Es hat eine Nebenbedingung je Konfiguration, also $q$ Stück.
Genau hier setzt der Beitrag von \textcite{gilmore_gomory1961,gilmore_gomory1963} an: Statt das LP
mit allen $q$ Spalten aufzustellen, löst man es über einer kleinen Spaltenteilmenge, liest die
Duallösung $y^\ast$ ab und sucht eine verletzte Nebenbedingung von \eqref{eq:gg-dual} durch
\[
  \max \Big\{ \textstyle\sum_{i=1}^{d} y^\ast_i\, a_i \;\Big|\;
  \sum_{i=1}^{d} p_i\, a_i \le C,\ \ 0 \le a_i \le u_i,\ \ a \in \mathbb{Z}^d \Big\},
\]
also durch ein beschränktes Rucksackproblem.
Ist der Optimalwert größer als $1$, so liefert das zugehörige $a$ eine neue Spalte und das
eingeschränkte LP wird erneut gelöst; ist er höchstens $1$, so ist $y^\ast$ für
\eqref{eq:gg-dual} zulässig und das eingeschränkte LP bereits optimal.
Dieses Verfahren heißt \emph{Spaltengenerierung}.

Für die vorliegende Arbeit wird es nicht benötigt.
Ist $d$ konstant, so lässt sich $\mathcal{K}(\Gamma, C)$ vollständig aufzählen: Wegen
$a_i \le \min(\lfloor C/p_i \rfloor, u_i)$ gilt
$q \le \prod_{i=1}^{d} \big(\min(\lfloor C/p_i \rfloor, u_i) + 1\big)$,
für festes $d$ also ein in $C$ polynomieller Wert.
Die Spaltengenerierung ist genau das Werkzeug, mit dem das Modell auch dann handhabbar bleibt,
wenn $d$ mit $n$ wächst und die Aufzählung ausscheidet.

\subsection{Eigenschaften des Konfigurations-LPs}
\label{sub:gg-eigenschaften}

\begin{lemma}[Schranken und Monotonie]
\label{lem:gg-schranken}
Für jede Instanz $\Gamma$ und jede Kapazität $C \ge p_{\max}$ gilt
\begin{enumerate}[label=(\alph*), nosep]
  \item $\big\lceil \mathrm{LP}[\Gamma, C] \big\rceil \le \mathrm{OPT}[\Gamma, C]$,
  \item $\mathrm{LP}[\Gamma, C] \ge \big(\sum_{j=1}^{n} p_j\big) / C$,
  \item $C \mapsto \mathrm{LP}[\Gamma, C]$ ist monoton fallend.
\end{enumerate}
\end{lemma}

\begin{proof}
(a) \eqref{eq:gg-lp} entsteht aus \eqref{eq:gg-ip} durch Weglassen der
Ganzzahligkeitsbedingung, also ist $\mathrm{LP}[\Gamma, C] \le \mathrm{OPT}[\Gamma, C]$;
da $\mathrm{OPT}[\Gamma, C]$ ganzzahlig ist, folgt die Aussage.

(b) Jede Konfiguration $a$ erfüllt $\sum_i a_i p_i \le C$.
Für eine zulässige Lösung $x$ folgt
$\sum_j p_j = \sum_i u_i p_i \le \sum_i p_i \sum_a a_i x_a
 = \sum_a x_a \sum_i a_i p_i \le C \sum_a x_a$.

(c) Für $C \le C'$ ist jede Konfiguration zur Kapazität $C$ auch eine zur Kapazität $C'$; die
Vielfachheitsschranken $a_i \le u_i$ hängen nicht von der Kapazität ab.
Der zulässige Bereich des LPs zu $C'$ enthält damit denjenigen zu $C$ (aufgefüllt mit Nullen),
und das Minimum kann nicht größer werden.
\end{proof}

Teil (b) besagt, dass das LP mindestens so scharf ist wie die Volumenschranke, aus der auch die
untere Intervallgrenze $c_{\mathrm{low}}^{(0)}$ in \eqref{eq:startintervall} stammt.
Teil (c) ist bemerkenswert im Licht von Bemerkung~\ref{bem:nicht-monoton}: Die \emph{Schranke}
$\mathrm{LP}[\Gamma, \cdot\,]$ fällt monoton, während $\mathrm{FFD}[\Gamma, \cdot\,]$ das nicht tut.
Für die gerundete Binzahl aus Abschnitt~\ref{sub:gg-rundung} wird hier keine Monotonie
behauptet; sie hängt von der jeweils gewählten Optimallösung ab.

Der eigentliche Hebel für die Rundung ist die geringe Zeilenzahl des LPs.

\begin{lemma}[Träger einer Optimalecke]
\label{lem:gg-ecke}
Das Konfigurations-LP besitzt eine optimale Ecke, und in jeder Ecke $x$ des zulässigen Bereichs
gilt $|\{a : x_a > 0\}| \le d$.
\end{lemma}

\begin{proof}
Der zulässige Bereich ist durch $d$ Bedarfszeilen und $q$ Vorzeichenbedingungen $x_a \ge 0$
beschrieben.
Eine \emph{Ecke} ist ein zulässiger Punkt, in dem $q$ linear unabhängige dieser Ungleichungen mit
Gleichheit erfüllt sind - gleichwertig: ein zulässiger Punkt, der sich nicht als echte
Konvexkombination zweier anderer zulässiger Punkte schreiben lässt.
Dass ein lösbares lineares Programm mit beschränkter Zielfunktion stets eine Optimallösung in einer
Ecke besitzt, ist Standard \cite[Kap.~11]{williamson_shmoys2011}.
Höchstens $d$ der $q$ straffen Ungleichungen können Bedarfszeilen sein; also sind mindestens
$q - d$ Vorzeichenbedingungen straff und ebenso viele Variablen gleich null.
\end{proof}

\noindent
Für das Konfigurations-LP des Bin Packing ist das die übliche Begründung des Rundungsverlusts;
\textcite[Abschnitt~4.6]{williamson_shmoys2011} verwenden dasselbe Argument.

\subsection{Rundung zu einer zulässigen Packung}
\label{sub:gg-rundung}

Aus einer optimalen Ecke $x^\ast$ wird in drei Schritten eine Packung gebaut.
Sei $F = \{a : x^\ast_a \notin \mathbb{Z}\}$ die Menge der fraktionalen Einträge; nach
Lemma~\ref{lem:gg-ecke} ist $|F| \le d$.

\begin{enumerate}[label=\textit{Schritt~\arabic*.}, leftmargin=*, itemsep=3pt]
  \item \emph{Abrunden.} Für jede Konfiguration $a$ werden $\lfloor x^\ast_a \rfloor$ Bins nach
    Muster $a$ gepackt. Weil \eqref{eq:gg-ip} in $\ge$-Form vorliegt, kann eine Größe dabei
    überdeckt werden; sind von Größe $p_i$ keine Aufgaben mehr übrig, so bleiben die
    entsprechenden Plätze des Musters unbesetzt.
  \item \emph{Restbedarf bestimmen.} Es sei
    $r_i = u_i - \sum_{a} \lfloor x^\ast_a \rfloor\, a_i$ für $r_i > 0$ und $r_i = 0$ sonst; das
    sind die nach Schritt~1 noch unplatzierten Aufgaben.
  \item \emph{Restbedarf packen.} Der Restbedarf wird auf zwei Arten gepackt - einmal mit einem
    Bin je fraktionaler Konfiguration $a \in F$, gefüllt nach deren Muster, und einmal mit FFD -
    und die kleinere der beiden Packungen wird verwendet.
    Bleibt ein Musterbin dabei leer, weil für seine Größen kein Restbedarf mehr besteht, so
    entfällt er; die Musterpackung braucht deshalb höchstens $|F|$ Bins, oft weniger.
\end{enumerate}

\noindent
Die Anzahl der insgesamt erzeugten Bins wird mit $\mathrm{GG}[\Gamma, C]$ bezeichnet.

\begin{lemma}[Rundungsschranke]
\label{lem:gg-rundung}
Schritt~3 packt den Restbedarf vollständig, und es gilt
\[
  \mathrm{OPT}[\Gamma, C] \;\le\; \mathrm{GG}[\Gamma, C] \;\le\;
  \big\lceil \mathrm{LP}[\Gamma, C] \big\rceil + |F| - 1 \;\le\;
  \mathrm{OPT}[\Gamma, C] + d - 1,
\]
sofern $F \neq \emptyset$; andernfalls ist
$\mathrm{GG}[\Gamma, C] = \mathrm{LP}[\Gamma, C] = \mathrm{OPT}[\Gamma, C]$.
\end{lemma}

\begin{proof}
Die untere Schranke ist klar, da das Verfahren eine zulässige Packung liefert.

\emph{Der Restbedarf passt in $|F|$ Bins.}
Aus der Zulässigkeit von $x^\ast$ folgt für jedes $i$
\[
  r_i \;\le\; u_i - \sum_{a} \lfloor x^\ast_a \rfloor\, a_i
  \;\le\; \sum_{a} \big(x^\ast_a - \lfloor x^\ast_a \rfloor\big)\, a_i
  \;=\; \sum_{a \in F} \big(x^\ast_a - \lfloor x^\ast_a \rfloor\big)\, a_i
  \;<\; \sum_{a \in F} a_i ,
\]
wobei die letzte Ungleichung strikt ist, falls die rechte Summe positiv ist.
Da beide Seiten ganzzahlig sind, gilt $r_i \le \sum_{a \in F} a_i$ für alle $i$.
Verteilt man den Restbedarf auf $|F|$ Bins, von denen der zu $a \in F$ gehörige höchstens $a_i$
Aufgaben der Größe $p_i$ aufnimmt, so wird jede verbliebene Aufgabe platziert und keine
Kapazität überschritten.
Die in Schritt~3 gewählte Packung ist höchstens so groß, also
$\mathrm{GG}[\Gamma, C] \le \sum_{a} \lfloor x^\ast_a \rfloor + |F|$.

\emph{Die obere Schranke.}
Für $a \notin F$ ist $\lceil x^\ast_a \rceil = \lfloor x^\ast_a \rfloor = x^\ast_a$, für $a \in F$
ist $\lceil x^\ast_a \rceil = \lfloor x^\ast_a \rfloor + 1 < x^\ast_a + 1$.
Summieren liefert
\[
  \sum_{a} \lfloor x^\ast_a \rfloor + |F| \;=\; \sum_{a} \lceil x^\ast_a \rceil
  \;<\; \mathrm{LP}[\Gamma, C] + |F| \;\le\; \lceil \mathrm{LP}[\Gamma, C] \rceil + |F| .
\]
Die linke Seite ist ganzzahlig und die erste Ungleichung strikt, also ist sie höchstens
$\lceil \mathrm{LP}[\Gamma, C] \rceil + |F| - 1$.
Mit Lemma~\ref{lem:gg-schranken}(a) und $|F| \le d$ folgt die letzte Abschätzung.
Ist $F = \emptyset$, so ist $x^\ast$ ganzzahlig, das Verfahren packt genau
$\mathrm{LP}[\Gamma, C]$ Bins, und mit Lemma~\ref{lem:gg-schranken}(a) ist dieser Wert optimal.
\end{proof}

\begin{bemerkung}[Warum nicht einfach aufrunden]
\label{bem:gg-aufrunden}
Die naheliegende Variante, jeden Eintrag aufzurunden und $\sum_a \lceil x^\ast_a \rceil$ Bins zu
packen, erfüllt dieselbe Schranke - die Rechnung im Beweis ist genau diese Summe.
Sie ist aber nie besser: Das Abrunden erzeugt $\sum_a \lfloor x^\ast_a \rfloor$ Bins, und der
Restbedarf braucht höchstens $|F|$ weitere, während das Aufrunden diese $|F|$ Bins immer belegt.
Deshalb wird hier abgerundet.
\end{bemerkung}

\subsection{Einsatz in MULTIFIT}
\label{sub:gg-multifit}

Als Subroutine erhält das Verfahren dieselbe Schnittstelle wie FFD und MFFD: Eingabe sind die
Aufgaben von $\Gamma$ und eine Kapazität $C$, Ausgabe ist eine Packung.
MULTIFIT akzeptiert $C$ genau dann, wenn $\mathrm{GG}[\Gamma, C] \le m$ gilt.
Sind die Aufgabengrößen ganzzahlig - wie auf den Instanzen aus Abschnitt~\ref{sub:hm-datensaetze} -,
so lässt eine Kapazität $C$ dieselben Packungen zu wie $\lfloor C \rfloor$; verwendet wird
dort deshalb die Variante der Binärsuche mit ganzzahligen Kapazitäten
(Abschnitt~\ref{sub:hm-aufbau}).

Gegenüber FFD ändert sich die Bedeutung einer Ablehnung.
Wegen $\lceil \mathrm{LP}[\Gamma, C] \rceil \le \mathrm{OPT}[\Gamma, C] \le \mathrm{GG}[\Gamma, C]$
zerfällt der Fall $\mathrm{GG}[\Gamma, C] > m$ in zwei Teile:
\begin{itemize}[leftmargin=*, itemsep=3pt]
  \item Ist bereits $\lceil \mathrm{LP}[\Gamma, C] \rceil > m$, so ist
    $\mathrm{OPT}[\Gamma, C] > m$, also nach Lemma~\ref{lem:dualitaet} $C < C^{*}_{\max}(\Gamma)$.
    Die Ablehnung ist dann \emph{zertifiziert}: Die untere Intervallgrenze wird auf einen Wert
    gesetzt, der nachweislich unterhalb des optimalen Makespans liegt.
    FFD zertifiziert eine Ablehnung erst ab $\mathrm{FFD}[\Gamma, C] > \tfrac{11}{9}\,m + \tfrac{6}{9}$:
    Nach \eqref{eq:ffd-guete} folgt daraus $\mathrm{OPT}[\Gamma, C] > m$.
    Im hier entscheidenden Fall $\mathrm{FFD}[\Gamma, C] = m+1$ folgt dagegen nichts
    über $\mathrm{OPT}[\Gamma, C]$.
  \item Ist $\lceil \mathrm{LP}[\Gamma, C] \rceil \le m < \mathrm{GG}[\Gamma, C]$, so bleibt offen,
    ob $C$ zulässig ist; die Rundung hat lediglich keine Packung mit $\le m$ Bins geliefert.
    Dieser Fall wird wie eine Ablehnung behandelt, damit die von MULTIFIT zurückgegebene Partition
    stets zur aktuellen oberen Intervallgrenze gehört.
\end{itemize}
Die Differenz $\mathrm{GG}[\Gamma, C] - \lceil \mathrm{LP}[\Gamma, C] \rceil$ misst damit genau
den Verlust, den die Rundung gegenüber der LP-Schranke verursacht.

Anders als bei FFD ist damit nicht gesichert, dass überhaupt eine Packung gespeichert wird:
Lemma~\ref{lem:startintervall}(b) gilt für FFD, und Lemma~\ref{lem:gg-rundung} lässt auch bei
$c_{\mathrm{high}}^{(0)}$ bis zu $m + d - 1$ Bins zu.
Wird jede getestete Kapazität abgelehnt, so gibt MULTIFIT keine Packung zurück; auf den
Instanzen aus Abschnitt~\ref{sub:hm-guete} tritt dieser Fall nicht auf.

Eine Gütegarantie für den Makespan folgt daraus nicht unmittelbar.
Die Makespan-Schranke $13/11 + 2^{-k}$ für MULTIFIT+FFD, die Theorem~\ref{thm:guete} mit
$r_m \le 13/11$ nach \textcite{cao1995} liefert, beruht auf
Implikation~\eqref{eq:ffd-akzeptiert}: Oberhalb von $r_m \cdot C^{*}_{\max}(\Gamma)$ kommt FFD
stets mit $m$ Bins aus, die Binärsuche akzeptiert dort also jede Kapazität und kann ihre obere
Intervallgrenze bis auf $r_m \cdot C^{*}_{\max}(\Gamma)$ absenken.
Damit sich diese Schlussweise übertragen lässt, müsste eine Konstante $r$ existieren, sodass
$C \ge r \cdot C^{*}_{\max}(\Gamma)$ stets $\mathrm{GG}[\Gamma, C] \le m$ nach sich zieht.
Für $r = 1$ liefert Lemma~\ref{lem:gg-rundung} nur $\mathrm{GG}[\Gamma, C] \le m + d - 1$ und
lässt damit gerade die $d-1$ Bins zu, die das Akzeptanzkriterium nicht mehr zulässt.
Größere Werte von $r$ sind damit jedoch nicht ausgeschlossen.
Sind $l_1, \ldots, l_m$ die Lasten einer optimalen Zuordnung, so trägt je ein Paar von ihnen
zusammen höchstens $2\,C^{*}_{\max}(\Gamma)$; es gilt also
$\mathrm{OPT}[\Gamma, C] \le \lceil m/2 \rceil$ für jedes $C \ge 2\,C^{*}_{\max}(\Gamma)$.
Mit Lemma~\ref{lem:gg-rundung} folgt $\mathrm{GG}[\Gamma, C] \le \lceil m/2 \rceil + d - 1$
und damit $\mathrm{GG}[\Gamma, C] \le m$, sobald $m \ge 2(d-1)$ gilt.
Für Instanzen mit wenigen verschiedenen Größen und hinreichend vielen Maschinen trägt die
Implikation also bereits mit $r = 2$.
Ein Gewinn ist das allerdings nicht: Nach Lemma~\ref{lem:untere-schranken} gilt
$c_{\mathrm{high}}^{(0)} = \sum_j p_j/m + p_{\max} \le 2\,C^{*}_{\max}(\Gamma)$, sodass schon
das Startintervall den Faktor~$2$ einhält.
Der Fall $r = 2$ zeigt daher nur, dass Lemma~\ref{lem:gg-rundung} die Schlussweise nicht
widerlegt; eine Schranke unterhalb von~$2$ liefert er nicht.
Ob eine von $d$ unabhängige Konstante $r < 2$ existiert, bleibt offen.
Der additive Term hängt von der Anzahl verschiedener Aufgabengrößen ab und ist nur für kleines $d$
klein - im Vergleich zur $\mathrm{OPT}+1$-Garantie von \textcite{jansen_solisoba2010} ist das
die schwächere Aussage, erkauft mit einem erheblich einfacheren Verfahren.
Wie oft die Schranke tatsächlich erreicht wird und ob der Rundungsverlust in der Praxis überhaupt
auftritt, ist eine empirische Frage.

\subsection{Implementierung}
\label{sub:gg-implementierung}

Die Implementierung liegt in \texttt{binpacking/gilmore\_gomory\_lp.py}.
Die Aufgabenliste wird zu $(p, u)$ aggregiert, $\mathcal{K}(\Gamma, C)$ per Tiefensuche mit
Kapazitätsschnitt aufgezählt - es werden also nur zulässige Vektoren erzeugt, statt die volle
Box aufzuspannen und nachträglich zu filtern - und das LP mit Gurobi gelöst.

Zwei Details sind für die Korrektheit wesentlich.
Erstens muss der Löser eine Ecke liefern.
Die Schranke $|F| \le d$ aus Lemma~\ref{lem:gg-ecke} gilt nur dort: Sind $x$ und $x'$ zwei
verschiedene optimale Ecken, so ist $\tfrac{1}{2}(x + x')$ ebenfalls optimal, aber keine Ecke und
kann bis zu $2d$ positive Einträge haben.
Simplexverfahren enden konstruktionsbedingt in einer Ecke \cite[Anh.~A]{williamson_shmoys2011}, ein Innere-Punkte-Verfahren ohne
anschließenden Crossover dagegen möglicherweise im Inneren einer optimalen Seitenfläche; Gurobi
wird deshalb auf die duale Simplex-Methode festgelegt.
Zweitens werden die $x^\ast_a$ mit einer Toleranz auf- und abgerundet, weil der Löser Werte wie
$2{,}9999999996$ statt $3$ zurückgibt.

Zusätzlich lässt sich mit denselben Spalten statt \eqref{eq:gg-lp} direkt \eqref{eq:gg-ip} lösen.
Das liefert $\mathrm{OPT}[\Gamma, C]$ exakt und damit eine Referenzlinie, gegen die sich sowohl
der Rundungsverlust des LPs als auch der Abstand von MULTIFIT zum bestmöglichen Ergebnis bei
jeder einzelnen getesteten Kapazität messen lässt.
Auf Instanzen mit kleinem $d$ ist das praktikabel.

Abschnitt~\ref{sec:hm} misst beides: den Rundungsverlust
$\mathrm{GG}[\Gamma,C] - \lceil\mathrm{LP}[\Gamma,C]\rceil$ und die Häufigkeit des
unentschiedenen Falls, dazu den Makespan gegen MULTIFIT+FFD und gegen die exakte
Referenzlinie.
